\section{Actual relative formal quotients by arbitrary original ideals}
\label{sec:hironaka-actual-relative-formal-quotients}

Write $B=\mathcal O_{\mathbb C^d,0}$, $R=\mathcal O_{\mathbb C^{d+t},0}$,
$\widehat B=\mathbb C[[u]]$, and $\widehat R=\mathbb C[[u,x]]$.
Let $T:R\to\widehat R$ be the actual Taylor homomorphism, let
$N=(x)$ be the original normal ideal, and let $\widehat N=(x)$ be
the same formal normal ideal. The $\widehat B$ action on
$\widehat R$ is the literal first-block renaming. Fix any whole
analytic ideal $I\subset R$; no finiteness, flatness, or chosen
normal form for $I$ is imposed.

\begin{theorem}[The whole formal ideal image is the actual coefficient span]
\label{hir:actual-relative-formal-whole-ideal-image}
\uses{hir:actual-formal-normal-truncation-coordinates,hir:actual-relative-polynomial-formal-taylor-coefficient}
\lean{Hironaka.HironakaActualRelativeFormalQuotients.formalTaylorIdealSpan_eq_wholeIdealImage}
\leanok
In $\widehat R/\widehat N^q$, the whole image of the extended ideal
$T(I)\widehat R$ is exactly the $\widehat B$-linear span of
the literal classes $[T(g)]$, $g\in I$.
\end{theorem}
\begin{proof}
Every displayed generator lies in the whole ideal image. That image
is stable under the original $\widehat B$ action, so it contains
the stated linear span $S$. Conversely $S$ is an ideal. To prove
closure under multiplication, first multiply a generating class
$[T(g)]$ by any class $[F]$. The whole formal normal-coordinate
surjectivity theorem writes
$[F]=\sum_{|\gamma|<q}c_\gamma[x^\gamma]$ with actual
$c_\gamma\in\widehat B$. Thus
\[
 [F][T(g)]=\sum_{|\gamma|<q}c_\gamma[T(x^\gamma g)].
\]
The analytic monomial $x^\gamma$ is the original normal monomial,
its Taylor series is the displayed formal monomial, and
$x^\gamma g\in I$. Every summand therefore belongs to $S$.
Distributivity and commutation of scalar multiplication with
multiplication extend this closure from the generating classes to
their entire $\widehat B$ span. Hence $S$ is an ideal containing
all classes $[T(g)]$ for $g\in I$. The ideal generated by those
classes, which is the whole image of $T(I)\widehat R$, is contained
in $S$. Both inclusions prove the equality, including $q=0$.
\end{proof}

\begin{theorem}[The actual ambient dictionary identifies the entire tensor ideal image]
\label{hir:actual-relative-formal-tensor-ideal-image}
\uses{hir:actual-formal-normal-truncation-base-change,hir:actual-relative-formal-whole-ideal-image}
\lean{Hironaka.HironakaActualRelativeFormalQuotients.map_tensorIdealRange_eq_wholeIdealImage}
\leanok
Put $A=R/N^q$ and let $I_q$ be the whole image of $I$ in $A$,
regarded as a $B$-submodule. The actual ambient linear dictionary
$E:\widehat B\otimes_B A\simeq_{\widehat B}\widehat R/\widehat N^q$
sends the entire image of $\widehat B\otimes_B I_q$ exactly to the
whole image of $T(I)\widehat R$.
\end{theorem}
\begin{proof}
A pure tensor $c\otimes a$ with $a\in I_q$ admits an actual
representative $g\in I$ with $a=[g]$: the original quotient map is
surjective, so its ideal image consists exactly of such images.
The actual ambient dictionary sends this tensor to $c[T(g)]$.
Every tensor is a finite sum of pure tensors, so the entire image
is contained in the $\widehat B$ span of the classes $[T(g)]$,
$g\in I$. Conversely every generating class $[T(g)]$ is the image
of $1\otimes[g]$ with $[g]\in I_q$. The image is a
$\widehat B$-submodule, so it contains that whole span. The preceding
whole-image theorem identifies this span with the asserted whole
formal ideal image, proving both inclusions.
\end{proof}

\begin{theorem}[Arbitrary original ideal quotients have actual relative formal base change]
\label{hir:actual-relative-formal-original-ideal-quotient}
\uses{hir:actual-relative-formal-tensor-ideal-image}
\lean{Hironaka.HironakaActualRelativeFormalQuotients.relativeFormalQuotientEquiv,Hironaka.HironakaActualRelativeFormalQuotients.relativeFormalQuotientEquiv_tmul_mk}
\leanok
For every original ideal $I$ and every $q\ge0$, there is an actual
$\widehat B$-linear isomorphism
\[
 \widehat B\otimes_B R/(I+N^q)\simeq_{\widehat B}
 \widehat R/(T(I)\widehat R+\widehat N^q),
 \qquad c\otimes[g]\longmapsto c[T(g)].
\]
\end{theorem}
\begin{proof}
Put $A=R/N^q$, and let $I_q$ be the whole image of $I$ in $A$.
The third isomorphism theorem identifies $A/I_q$ with
$R/(I+N^q)$, preserving every original numerator. Right exact
module tensor quotient identifies $\widehat B\otimes_B(A/I_q)$
with the quotient of $\widehat B\otimes_B A$ by the entire image
of $\widehat B\otimes_B I_q$.

The preceding actual image theorem proves that the unchanged ambient
linear dictionary takes this submodule exactly onto the whole image
of $T(I)\widehat R$ in $\widehat R/\widehat N^q$. A linear
isomorphism descends through submodules with these equal images,
so it gives an isomorphism of the two entire quotient modules.
The third isomorphism theorem identifies the latter quotient with
$\widehat R/(T(I)\widehat R+\widehat N^q)$.
Composing gives the asserted isomorphism. Every quotient and tensor
map retains its representative, while the ambient dictionary sends
$c\otimes[g]$ to $c[T(g)]$; hence the composite has exactly the
displayed original-numerator formula.
\end{proof}

\begin{theorem}[All normal-power transition maps commute]
\label{hir:actual-relative-formal-normal-transitions}
\uses{hir:actual-relative-formal-original-ideal-quotient}
\lean{Hironaka.HironakaActualRelativeFormalQuotients.relativeFormalQuotientEquiv_transition}
\leanok
For $s\le q$, the order-$q$ to order-$s$ quotient maps on both sides
commute with the same original-ideal relative formal dictionaries.
\end{theorem}
\begin{proof}
The inclusions $N^q\subset N^s$ and
$\widehat N^q\subset\widehat N^s$ define the actual quotient maps,
with the same unchanged original ideal on each side. On a pure
tensor $c\otimes[g]$, applying the analytic transition first and
then the order-$s$ dictionary gives $c[T(g)]$ in the order-$s$
target. Applying the order-$q$ dictionary first and then the formal
transition gives exactly the same class. Both composites are
$\widehat B$-linear. Since every tensor is a finite sum of pure
tensors, additivity gives equality on every tensor, proving the
whole transition identity.
\end{proof}

\begin{remark}
These are whole quotient and transition dictionaries for arbitrary
original analytic ideals. They do not select a canonical standard
normal-monomial family or prove the reduced-vertex centre-power
condition. Those conclusions require an additional comparison of
the actual chosen monomial families and reduced-series support.
\end{remark}
